\chapter[Termo de Abertura do Projeto]{Termo de Abertura do Projeto}
\vspace{-1cm}

\section{Dados do Projeto}
\begin{description}
    \item [Nome do Projeto:] RATOBÔ
    \item [Data de Abertura:] 11/05/2026
    \item [Código:] 1-A
    \item [Patrocinador:] Universidade de Brasília (UnB)
    \item [Gerente:] Gabriel Alves de Araujo / Matrícula: 241025523 / E-mail: 241025523@aluno.unb.br / Contato: (61) 98469-7666
\end{description}

\section{Objetivos}
O objetivo principal é desenvolver um robô \textit{Micromouse} autônomo para mapear e resolver labirintos, integrado a uma aplicação web de telemetria em tempo real, sob gestão via Estrutura Analítica do Projeto (EAP).

Para fins de rastreabilidade, os objetivos específicos seguem a metodologia SMART:

\begin{description}
    \item [\textit{Specific} (Específico):] Criar um robô autônomo (máx. 25 cm) para navegar em malhas $4\times4$, $8\times8$ e $16\times16$ sem danificar a pista. Inclui interface web e banco de dados para telemetria instantânea (coordenadas, energia, velocidade, tempo e status).
    \item [\textit{Measurable} (Mensurável):] Resolução autônoma dos três labirintos, integridade dos registros de "Desafio" e "Trajeto" no banco de dados e construção de uma pista de testes física $4\times4$ (chão preto, paredes brancas de 5 cm, topos vermelhos).
    \item [\textit{Agreed} (Acordado):] Escopo e requisitos alinhados entre a equipe multidisciplinar e as competências exigidas pela disciplina de Projeto Integrador 1 (PI1).
    \item [\textit{Realistic} (Realista):] Viabilidade financeira assegurada pelo uso de componentes de baixo custo. A complexidade algorítmica é compatível com o nível técnico da equipe e com o cronograma de desenvolvimento.
    \item [\textit{Time Bound} (Limitado no Tempo):] Entregas incrementais vinculadas ao calendário acadêmico, garantindo protótipo e testes operacionais para a avaliação e entrega final na data estipulada pela disciplina.
\end{description}

\section{Mercado-Alvo}
O ecossistema foca em estudantes, professores e pesquisadores da área tecnológica:
\begin{itemize}
    \item \textbf{Estudantes (Engenharias):} Ferramenta prática de aprendizado interdisciplinar e validação de sistemas de controle.
    \item \textbf{Instituições de Ensino:} Plataforma modular \textit{open-source} para uso como recurso didático em laboratórios de automação.
    \item \textbf{Pesquisadores:} Base de hardware e simulador web para experimentação e avaliação comparativa de algoritmos de inteligência computacional.
\end{itemize}

\section{Requisitos do Sistema}
A Tabela \ref{tab:requisitos_micromouse} consolida as necessidades de negócio e restrições técnicas mapeadas para o projeto.

\begin{table}[H] 
    \centering
    \small
    \caption{Requisitos funcionais (RF) e não funcionais (RNF) do sistema Micromouse.}
    \label{tab:requisitos_micromouse}
    \renewcommand{\arraystretch}{1.3}
    \begin{tabularx}{\linewidth}{|l|l|X|}
        \hline
        \textbf{Cód.} & \textbf{Categoria} & \textbf{Descrição do Requisito} \\ \hline
        RF-01 & Navegação Autônoma & Resolver labirintos ($4\times4$, $8\times8$, $16\times16$) autonomamente, mapeando paredes e detectando o centro. \\ \hline
        RF-02 & Interface Web & Exibir telemetria em tempo real: tipo de labirinto, trajeto, bateria, velocidade, tempo e status de conclusão. \\ \hline
        RF-03 & Banco de Dados & Armazenar os dados finais de cada execução ao término do desafio. \\ \hline
        RF-04 & Histórico & Permitir consulta na aplicação web detalhando corridas armazenadas. \\ \hline
        RNF-01 & Dimensões & O robô não deve exceder 25 cm de comprimento/largura (sem restrição de altura). \\ \hline
        RNF-02 & Locomoção & Proibido voar, pular, escalar, usar combustão ou danificar o labirinto. \\ \hline
        RNF-03 & Inalterabilidade & Proibido alterar código/memória durante a corrida (exceto em falhas críticas). \\ \hline
        RNF-04 & Pista Física & Construir pista $4\times4$ (chão preto, paredes brancas de 5 cm com topos vermelhos). \\ \hline
        RNF-05 & Autoria & Desenvolvimento integral pela equipe (proibido plágio ou soluções prontas). \\ \hline
        RNF-06 & Gestão (GitHub) & Versionamento e controle de tarefas via \textit{issues} (máx. 2 membros por \textit{issue}). \\ \hline
    \end{tabularx}
\end{table}

\section{Justificativa}
O projeto consolida de forma prática os conceitos multidisciplinares da FCTE, unindo Engenharia de Software (arquitetura web e conectividade), Eletrônica (microcontroladores), Energia (eficiência) e Aeroespacial (cinemática). Além disso, preenche uma lacuna didática ao oferecer uma plataforma \textit{open-source} de baixo custo que, diferentemente de kits comerciais restritos, exige o domínio da construção do protótipo, gerando um ecossistema documentado para pesquisas em robótica universitária.

\section{Indicadores de Sucesso}
Para fins de avaliação técnica e acadêmica, estabelecem-se os seguintes indicadores:
\begin{enumerate}
    \item Índice de assertividade do mapeamento espacial nas malhas $4\times4$, $8\times8$ e $16\times16$.
    \item Tempo médio de resposta do \textit{pipeline} de telemetria (\textit{delay} WebSocket).
    \item Taxa de integridade dos pacotes persistidos no banco de dados PostgreSQL.
    \item Quantidade de colisões físicas registradas por ciclo de corrida.
    \item Erro percentual da odometria e sensores frente às dimensões reais da pista.
    \item Eficiência energética (Ah) em função da velocidade média dos motores.
    \item Índice de cobertura de testes automatizados (\textit{coverage}) alcançado.
    \item Vazão de tarefas (\textit{throughput}) medida pelo fechamento de \textit{issues} no GitHub.
\end{enumerate}